\documentclass[11pt]{article}
\usepackage[letterpaper,margin=0.82in]{geometry}
\usepackage{amsmath,amssymb,amsthm,mathtools}
\usepackage{graphicx}
\usepackage{microtype}
\usepackage[hidelinks]{hyperref}
\newtheorem{theorem}{Theorem}
\newtheorem{lemma}[theorem]{Lemma}
\newcommand{\Om}{\Omega_n}
\newcommand{\M}{\mathcal M_n}
\newcommand{\OO}{\mathcal O}
\newcommand{\C}{\mathbb C}
\newcommand{\GL}{\mathrm{GL}_n(\C)}
\newcommand{\Aut}{\operatorname{Aut}}
\newcommand{\Pic}{\operatorname{Pic}}
\newcommand{\Spec}{\operatorname{Spec}}
\title{Normalized spectral-ball automorphisms are\\global holomorphic conjugations}
\author{Candidate solution packet for arXiv:1011.2485}
\date{13 August 2026}
\begin{document}
\maketitle

\noindent\textbf{Status.}
Candidate full solution, likely valid, human review needed.

\section{Source question and scope}
Let
\[
 \Om=\{X\in\M(\C):\rho(X)<1\}
\]
be the spectral ball.  In the final paragraph of Kosi\'nski's note
\cite{Kosinski2010} (PDF page 3), the first of two problems asks:
\begin{quote}
Is every normalized automorphism of the spectral ball of the form
$X\mapsto u(X)Xu(X)^{-1}$, where $u:\Om\to\GL$ is holomorphic?
\end{quote}
Here normalized means $F(0)=0$ and $F'(0)=\mathrm{id}$.  The source evidence
is reproduced below.

\begin{center}
\includegraphics[width=0.91\textwidth]{figures/open_problem_crop.png}
\end{center}

The theorem below answers this first problem affirmatively for every $n$.
It does not classify all holomorphic $u$ whose conjugation map is an
automorphism, which is the source's separate second problem.

\section{Proof intuition}
On the open set of cyclic matrices, conjugators exist locally and any two
differ by an invertible matrix commuting with $X$.  The commuting algebra of a
cyclic matrix is exactly the polynomial algebra generated by $X$.  Globally,
this is the function algebra of the spectral cover
\[
 Y=\{(X,\lambda):\det(\lambda I-X)=0\}\longrightarrow\Om.
\]
Thus the failure of the local conjugators to glue is measured by a line
bundle on the part of $Y$ lying over cyclic matrices.

Two ingredients kill that obstruction.  First, the normalized rescaling
$F_t(X)=t^{-1}F(tX)$ deforms $F$ through spectrum-preserving automorphisms to
the identity; local conjugators chosen by Krylov bases make the obstruction
line bundle vary continuously, so its first Chern class is zero.  Second, the
deleted part of the spectral cover has codimension three.  A depth/local
cohomology argument gives $H^1(\OO)=0$ on its complement, making the Chern
class map on the Picard group injective.  The obstruction is therefore
holomorphically trivial.  After gluing on cyclic matrices, ordinary Hartogs
extension crosses the codimension-three derogatory locus; purity of the zero
set of a holomorphic function ensures that the extended conjugator remains
invertible.

\section{Geometric lemmas}
Call $X$ \emph{cyclic} if some $v\in\C^n$ makes
$v,Xv,\ldots,X^{n-1}v$ a basis, and let $D\subset\Om$ be the derogatory
(noncyclic) locus.  Put $U=\Om\setminus D$.

\begin{lemma}[Codimension and the spectral cover]\label{lem:geometry}
For $n\ge2$, $D$ is analytic of codimension three.  The finite flat spectral
cover
\[
 \pi:Y:=\{(X,\lambda)\in\Om\times\C:
          \det(\lambda I-X)=0\}\longrightarrow\Om
\]
has degree $n$.  It is a Stein Cohen--Macaulay, normal hypersurface.  If
$W=\pi^{-1}(D)$ and $V=Y\setminus W=\pi^{-1}(U)$, then
\[
 \operatorname{codim}_Y W=3,
 \qquad H^1(V,\OO_V)=0.
\]
\end{lemma}

\begin{proof}
A matrix is derogatory precisely when some eigenvalue has geometric
multiplicity at least two.  Hence $D$ is the projection of
\[
 S=\{(X,\lambda):\operatorname{rank}(X-\lambda I)\le n-2\}.
\]
For fixed $\lambda$, the rank condition has codimension four in $\M$.
Consequently $S$ has dimension $n^2-3$ in $\M\times\C$.  Its projection has
finite fibers, so $D$ has dimension $n^2-3$ (the generic stratum with one
diagonalizable double eigenvalue realizes this dimension).

The equation defining $Y$ is monic of degree $n$ in $\lambda$, so $\pi$ is
finite flat of degree $n$.  The spectral ball is pseudoconvex and therefore
Stein; hence $Y$, a closed analytic hypersurface in $\Om\times\C$, is Stein.
Every hypersurface is Cohen--Macaulay.  Differentiating the determinant shows
that a point of $Y$ is singular exactly when
$\operatorname{rank}(\lambda I-X)\le n-2$: the matrix derivatives are the
entries of the adjugate.  Thus the singular locus has codimension three in
$Y$, and Serre's hypersurface criterion implies that $Y$ is normal.

Because $\pi$ is finite and $D$ has codimension three, so does
$W=\pi^{-1}(D)$.  The analytic Cohen--Macaulay depth theorem
\cite{BanicaStanasila} gives
\[
 H^q_W(Y,\OO_Y)=0\qquad(q<3).
\]
Cartan's theorem B \cite{GrauertRemmert} gives
$H^q(Y,\OO_Y)=0$ for $q>0$.  The long exact
local-cohomology sequence for $(Y,W)$ now yields
\[
 H^1(V,\OO_V)=0,
\]
as asserted.  Notice that for $n>2$, $W$ can strictly contain
$\operatorname{Sing}Y$; only its codimension and the Cohen--Macaulay depth are
used.
\end{proof}

\begin{lemma}[Centralizers and the obstruction bundle]\label{lem:bundle}
Let $G:U\to\M$ be holomorphic, spectrum-preserving, and pointwise conjugate to
the identity map.  Local holomorphic conjugators
$G(X)=p_i(X)Xp_i(X)^{-1}$ determine a holomorphic line bundle $L_G$ on $V$.
The bundle is trivial if and only if there is one holomorphic
$p:U\to\GL$ satisfying $G(X)=p(X)Xp(X)^{-1}$.
\end{lemma}

\begin{proof}
Local conjugators exist directly.  At $X_0\in U$, choose a vector $v$ cyclic
for both $X_0$ and $G(X_0)$; such a vector exists because each set of cyclic
vectors is a nonempty Zariski-open subset of $\C^n$.  After shrinking, the
Krylov matrices
\[
 K_v(A)=[v,Av,\ldots,A^{n-1}v]
\]
are invertible for $A=X$ and $A=G(X)$.  Since $X$ and $G(X)$ have the same
characteristic polynomial,
\[
 p(X)=K_v(G(X))K_v(X)^{-1}
\]
satisfies $G(X)=p(X)Xp(X)^{-1}$.

On an overlap, $c_{ij}=p_i^{-1}p_j$ commutes with $X$.  For cyclic $X$, its
centralizer is $\C[X]$.  In sheaf form this gives the canonical algebra
isomorphism
\[
 \pi_*\OO_V\ \cong\
 \OO_U[t]/(\det(tI-X))\ \cong\
 \{A\in\M(\OO_U):AX=XA\}.
\]
Thus every $c_{ij}$ is a unit in $\pi_*\OO_V$, equivalently a nowhere-zero
holomorphic transition function on $V$.  The cocycle relation defines a line
bundle $L_G$.

If $L_G$ is trivial, write $c_{ij}=h_i^{-1}h_j$ with $h_i$ units of
$\pi_*\OO_V$, hence invertible centralizer matrices.  Then
$p_i h_i^{-1}=p_jh_j^{-1}$ on overlaps, and these matrices glue to the desired
global conjugator.  The converse follows by comparing any local conjugators
with a global one.
\end{proof}

\section{Main theorem}
\begin{theorem}\label{thm:main}
For every $n\ge1$, every normalized holomorphic automorphism
$F\in\Aut(\Om)$ admits a holomorphic map $u:\Om\to\GL$ such that
\[
 F(X)=u(X)Xu(X)^{-1}\qquad(X\in\Om).
\]
\end{theorem}

\begin{proof}
The case $n=1$ is immediate.  Assume $n\ge2$.  A normalized automorphism of
the spectral ball preserves the spectrum \cite{RansfordWhite}; in particular,
on $U$ it sends every cyclic $X$ to a cyclic matrix conjugate to $X$.

For $0<t\le1$ define
\[
 F_t(X)=t^{-1}F(tX),
 \qquad F_0(X)=X.
\]
The Taylor expansion at zero and $F'(0)=\mathrm{id}$ show continuity at
$t=0$.  Spectrum preservation gives
$\operatorname{Spec}F_t(X)=\operatorname{Spec}X$.  Moreover
$t^{-1}F^{-1}(tX)$ is the inverse of $F_t$, so every $F_t$ is an automorphism.

Apply Lemma~\ref{lem:bundle} to $F_t|_U$.  The local Krylov construction in
its proof is parametric: near any $(t_0,X_0)\in[0,1]\times U$, choose one
vector cyclic for both $X_0$ and $F_{t_0}(X_0)$; invertibility persists nearby.
The resulting transition functions therefore define a continuous complex
line bundle $\mathcal L$ on $[0,1]\times V$, whose restriction at $t$ is the
underlying topological bundle of $L_{F_t}$.  At $t=0$, the Krylov formula gives
the identity conjugator, so $\mathcal L|_{\{0\}\times V}$ is trivial.
Homotopy invariance now gives
\[
 c_1(L_F)=c_1(L_{F_1})=0\quad\text{in }H^2(V,\mathbb Z).
\]

The exponential sequence on $V$ contains
\[
 H^1(V,\OO_V)\longrightarrow \Pic(V)
   \xrightarrow{c_1} H^2(V,\mathbb Z).
\]
Lemma~\ref{lem:geometry} makes the left group zero, so $c_1$ is injective.
It follows that $L_F$ is holomorphically trivial.  Lemma~\ref{lem:bundle}
therefore supplies a holomorphic $u:U\to\GL$ with
$F(X)=u(X)Xu(X)^{-1}$.

Finally, $D=\Om\setminus U$ has codimension three.  Hartogs extension
\cite{GrauertRemmert} extends
every matrix entry of $u$ holomorphically across $D$, giving
$\widetilde u:\Om\to\M$.  Its determinant is nonzero on $U$.  If
$\det\widetilde u$ vanished, its zero set would be a divisor, but it would
have to lie inside the codimension-three set $D$, which is impossible.
Therefore $\widetilde u$ takes values in $\GL$, and the conjugation identity
extends to all of $\Om$ by holomorphic continuity.
\end{proof}

\section{Verification, limitations, and novelty}
The proof was stress-tested in all dimensions, at ramified cyclic fibers, and
at the derogatory locus.  The distinction
$W=\pi^{-1}(D)\supseteq\operatorname{Sing}Y$ for $n>2$ is explicit; the depth
argument applies to $W$ itself.  No computation is used as proof.

A bounded literature search on 13 August 2026 checked exact and close variants
of ``spectral ball global holomorphic conjugation'', the source
\cite{Kosinski2010}, Thomas's local theorem \cite{Thomas2008}, and Kosi\'nski's
later $n=2$ local/density theorem \cite{Kosinski2012}.  The latter states that
automorphisms of $\Omega_2$ are \emph{locally} holomorphic conjugations and
that certain conjugations generate a dense subgroup; it does not state the
global theorem above.  No exact prior answer was found.  Novelty confidence is
moderate because the search was bounded.

The source's second final problem, a classification of all holomorphic $u$
whose conjugation map is an automorphism, remains outside this packet.  Human
review should concentrate on the local-cohomology vanishing in
Lemma~\ref{lem:geometry} and the parametric line-bundle construction in
Theorem~\ref{thm:main}.

\begin{thebibliography}{9}
\bibitem{Kosinski2010}
L.~Kosi\'nski, \emph{Note on the group of automorphisms of the Spectral Ball},
arXiv:1011.2485 (2010).
\bibitem{RansfordWhite}
T.~J. Ransford and M.~C. White, \emph{Holomorphic self-maps of the spectral
unit ball}, Bull. London Math. Soc. 23 (1991), 256--262.
\bibitem{Thomas2008}
P.~J. Thomas, \emph{A local form for the automorphisms of the spectral unit
ball}, Collect. Math. 59 (2008), 321--324; arXiv:0801.3396.
\bibitem{Kosinski2012}
L.~Kosi\'nski, \emph{A local form of automorphisms of the spectral unit ball},
arXiv:1202.5793 (2012).
\bibitem{BanicaStanasila}
C.~B\u anic\u a and O.~St\u an\u a\c sil\u a,
\emph{Algebraic Methods in the Global Theory of Complex Spaces},
Wiley, 1976.
\bibitem{GrauertRemmert}
H.~Grauert and R.~Remmert, \emph{Coherent Analytic Sheaves}, Springer, 1984.
\end{thebibliography}
\end{document}
